% from ega4-10-fr.tex
\oldpage[IV-3]{105}
fini de ses points fermés et de son point générique, donc \emph{f)} implique
\emph{e)} sans supposer $A$ noethérien. Enfin \emph{e)} implique \emph{d)} sans
supposer $A$ noethérien, car le point générique $x$ de $X$ est le complémentaire
de la réunion des adhérences $\overline{\{y\}}$, où $y$ parcourt l'ensemble des
points $y\neq x$, et comme ces points sont en nombre fini, $\{x\}$ est
complémentaire d'un ensemble fermé dans $X$.

\begin{corollary}[10.5.2]
\label{IV.10.5.2-fr}
Soit $A$ un anneau noethérien. Pour que $A$ soit un anneau de Jacobson, il faut
et il suffit qu'il n'existe aucun idéal premier $\mathfrak p$ de $A$ tel que
$A/\mathfrak p$ soit un anneau semi-local de dimension $1$.
\end{corollary}

Cela résulte aussitôt de (10.5.1) et de la condition \emph{b)} de (10.4.5),
les idéaux premiers $\mathfrak p$ de $A$ tels que $A/\mathfrak p$ soit
semi-local de dimension $0$ étant les idéaux maximaux de $A$.

\begin{corollary}[10.5.3]
\label{IV.10.5.3-fr}
Soit $X$ un préschéma noethérien irréductible. Les conditions suivantes sont
équivalentes~:
\begin{enumerate}[label=\emph{\alph*)}]
  \item Le point générique de $X$ est isolé.
  \item $X$ est fini.
  \item $X$ est de dimension $\leq1$ et l'ensemble de ses points fermés est
  fini.
\end{enumerate}
\end{corollary}

Il existe par hypothèse un nombre fini d'ouverts affines irréductibles $U_i$
($1\leq i\leq n$) recouvrant $X$, et dont chacun contient donc le point
générique de $X$~; il suffit de prouver l'équivalence de \emph{a)}, \emph{b)}
et \emph{c)} pour chacun des $(U_i)_{\mathrm{red}}$ (compte tenu de
($\mathbf 0$, 14.1.7)). Mais cette équivalence résulte alors de (10.5.1).

\begin{remark}[10.5.4]
\label{IV.10.5.4-fr}
Un préschéma noethérien $X$ vérifiant les conditions équivalentes de (10.5.3)
n'est pas nécessairement un schéma affine~; en fait il peut même être
\emph{non séparé}. On en a un exemple en remplaçant, dans l'exemple
(\textbf{I}, 5.5.11) de «~droite affine avec point dédoublé~», $X_1$ et $X_2$
par le spectre de l'anneau de valuation discrète $(K[s])_{(s)}$, $U_{12}$ et
$U_{21}$ par l'ouvert réduit au point générique dans $X_1$ et $X_2$
respectivement~; le préschéma non séparé $X$ que l'on obtient a exactement
$3$ points.
\end{remark}

\begin{proposition}[10.5.5]
\label{IV.10.5.5-fr}
Soit $X$ un préschéma noethérien.
\begin{enumerate}[label=(\roman*)]
  \item L'ensemble $X_0$ des points $x\in X$ tels que $\overline{\{x\}}$ soit
  fini est très dense dans $X$.
  \item Pour que $X$ soit un préschéma de Jacobson, il faut et il suffit qu'il
  n'existe aucun sous-préschéma de $X$ isomorphe au spectre d'un anneau
  semi-local intègre de dimension $1$.
\end{enumerate}
\end{proposition}

La condition que $\overline{\{x\}}$ soit fini équivaut en effet ici (10.5.3)
au fait que $x$ soit isolé dans $\overline{\{x\}}$, $\overline{\{x\}}$ étant
espace sous-jacent d'un sous-préschéma (noethérien) de $X$~; l'assertion (i)
résulte donc de (5.1.10.1). De même, compte tenu de (10.4.5.1), pour démontrer
l'assertion (ii), remarquons que pour qu'un point non fermé $x$ de $X$
appartienne à $X_0$, il faut et il suffit que le sous-préschéma fermé intègre
$Y$ de $X$ ayant $\overline{\{x\}}$ pour espace sous-jacent soit de dimension
$1$ et fini, donc réunion finie de sous-préschémas ouverts (dans $Y$) affines
$U_i$ qui sont des spectres d'anneaux semi-locaux intègres de dimension $1$.
Inversement, s'il y a un sous-préschéma $Z$ de $X$ qui soit spectre d'un anneau
semi-local intègre de dimension $1$, $Z$ n'est pas un préschéma de Jacobson
(10.5.2), donc il en est de même de $X$ (10.3.2).

